%% kapa: exam=1 round=35 year=2024 period=1 form=A subject=real_estate subject_ko=부동산학원론 seq=9 number=89
%% source: Q-Net 감정평가사 기출문제 (cst00309) artlSeq=5213229 file=raw/1008/5213229/1교시 문제지 원본.pdf
%% conversion: kordoc 4.24.51 -> tools/gen1.py (auto); review=proofread-sonnet; flags=none
\begin{question}{89}
외부효과에 관한 내용으로 ( )에 들어갈 것으로 옳은 것은?
\begin{framedtext}
○ 부동산의 특성 중에서 ( ㄱ )은 외부효과를 발생시킨다. \par
○ 부동산시장 참여자가 자신들의 행동이 초래하는 외부효과를 의사결정에서 감안하도록 만드는 과정을 외부효과의 ( ㄴ )라 한다. \par
\end{framedtext}
\begin{choices}
  \item ㄱ: 부동성, ㄴ: 유동화
  \item ㄱ: 부동성, ㄴ: 내부화
  \item ㄱ: 인접성, ㄴ: 유동화
  \item ㄱ: 개별성, ㄴ: 내부화
  \item ㄱ: 개별성, ㄴ: 유동화
\end{choices}
\end{question}
\answer{2}
